
\documentclass[12pt,a4paper]{report}

% ==== PAQUETES BÁSICOS ====

\usepackage{subcaption}
\usepackage[utf8]{inputenc}
\usepackage{amsmath,amssymb}
\usepackage[T1]{fontenc}
\usepackage[spanish]{babel}
\usepackage{graphicx} % para imágenes
\usepackage{caption}
\usepackage{hyperref} % hipervínculos
\usepackage{geometry} % márgenes
\usepackage{booktabs}
\usepackage{subcaption}
\geometry{left=3cm,right=2.5cm,top=3cm,bottom=2.5cm}

% ==== CÓDIGO FUENTE ====
\usepackage[linesnumbered,ruled,vlined]{algorithm2e}
\SetKwInOut{KwIn}{Entrada}
\SetKwInOut{KwOut}{Salida}

\usepackage{listings}
\usepackage{xcolor}
\usepackage{tikz}
%%%%%Uso de bonificador de javascript
\usepackage{listings}
\usepackage{xcolor}
\lstdefinestyle{javascript}{
  language=Java,
  morekeywords={async,await,import,export,from},
  keywordstyle=\color{blue}\bfseries,
  stringstyle=\color{red},
  commentstyle=\color{gray},
  basicstyle=\ttfamily\small,
  breaklines=true,
  showstringspaces=false
}
%%%%%%%%%%%%%%%%%%%%%%%%%%%%%%%%


\usetikzlibrary{shapes.geometric, arrows,positioning}
\tikzstyle{startstop} = [rectangle, rounded corners, minimum width=2cm, minimum height=0.7cm,text centered, draw=black, fill=red!30]
\tikzstyle{process} = [rectangle, minimum width=2cm, minimum height=0.7cm, text centered, draw=black, fill=blue!20]
\tikzstyle{decision} = [diamond, minimum width=3cm, minimum height=0.7cm, text centered, draw=black, fill=green!30]
\tikzstyle{arrow} = [thick,->,>=stealth]
\lstset{
    language=C,                 % Puedes cambiarlo: Python, Java, etc.
    basicstyle=\ttfamily\small,   % Fuente monoespaciada
    keywordstyle=\color{blue},    % Palabras clave
    stringstyle=\color{red},      % Strings
    commentstyle=\color{gray},    % Comentarios
    backgroundcolor=\color{black!5},
    numbers=left,                 % Números de línea
    numberstyle=\tiny\color{gray},
    stepnumber=1,
    breaklines=true,
    tabsize=4
}
\begin{document}
% ==== PORTADA ====
\begin{titlepage}
    \centering
    
    % Imagen institucional (ejemplo: logo)
    \includegraphics[width=0.3\textwidth]{ug.png}\par\vspace{1cm}
    
    % Nombre de la universidad
    {\Large \textbf{Universidad de Guanajuato}}\\[0.5cm]
    
    % Nombre de la materia
    {\large \textbf{Materia: Aprendizaje de máquina}}\\[1cm]
    {\large \textbf{Carrera: Licenciatura en Ingeniería de Datos e Inteligencia Artificial}}\\[3cm]
    
    % Título de la tarea
    {\Huge \textbf{Tarea 2: Uso de aprendizaje no supervisado para identificación de tareas}}\\[2cm]
    
    % Datos del alumno
    {\large \textbf{Alumno:} José Adrián Rodríguez González}\\[0.3cm]
    {\large \textbf{NUA:} 148661}\\[1.5cm]
    
    % Datos del profesor
    {\large \textbf{Profesor:} Dr. Dora Almanza Luz Ojeda}\\[3cm]
    
    % Fecha
    {\large \today}
    
\vfill
\end{titlepage}

\tableofcontents % Índice automático
\newpage

% ==== INTRODUCCIÓN ====
\chapter{Introducción y objetivos}
\end{document}
